
\SimpleSystemTeX 允许用户通过在导言区使用 \cs{SetStyle} 命令自定义\newconcept{样式代码}{Style Code}, 从而设定或修改文档各组件的样式. 除了标准\LaTeX 命令之外, 用户还可以在样式代码中使用以 \texttt{@} 为前缀的\newconcept{显示参数}{Display Argument}来动态地提取文档内容, 这些参数的值同样能用于条件判断, 以实现高度灵活的排版效果.

\UserCommand{设定样式}<\cs{SetStyle}>{
    \cs{SetStyle}
    \{\meta{样式名称}\}
    \{\meta{样式代码}\longarg\}
}[
    - \meta{样式名称} : 需要设定或修改的样式名称.
][
    - \meta{样式代码}\longarg\ : 用于定义该样式具体外观的\LaTeX 代码. 您可以在其中使用用于排版的标准\LaTeX 命令, 带有 \texttt{@} 前缀的显示参数以及下列带有 \texttt{\#} 前缀的条件判断语句:
        - - \texttt{\#ifempty} : 判断显示参数是否为空.
        \LocalUsage{ifempty}(\#ifempty){
            \upshape
            \#ifempty
            @\meta{显示参数}
            \{\meta{真分支代码}\longarg\}
            \{\meta{假分支代码}\longarg\}
        }
        - - \texttt{\#ifis} : 判断显示参数展开一次后是否等于某值.
        \LocalUsage{ifis}(\#ifis){
            \upshape
            \#ifis
            @\meta{显示参数}
            \{\meta{值}\}
            \{\meta{真分支代码}\longarg\}
            \{\meta{假分支代码}\longarg\}
        }
        - - \texttt{\#ifequal} : 判断两个显示参数展开一次后是否相等.
        \LocalUsage{ifequal}(\#ifequal){
            \upshape
            \#ifequal
            @\meta{显示参数1\,}
            @\meta{显示参数2\,}
            \{\meta{真分支代码}\longarg\}
            \{\meta{假分支代码}\longarg\}
        }
]

在编写样式代码时, 请务必保证显示参数不紧跟在任何宏命令后面, 否则可能会导致语法解析错误. 您需要在 \texttt{@} 前添加一个空格来避免这种情况, 例如 \cs{par}\texttt{ @\meta{显示参数}} .

\newpage 可供设定的样式名称及其可用的显示参数包括:

\List{目录样式}(\large 目录){
    \par\medskip{}
    - \textbf{\texttt{TOCTitle} : 目录标题}.
    - \textbf{\texttt{TOCPartTitle} : 目录中的分块标题}.
        - - \texttt{@part} : 分块.
    - \textbf{\texttt{TOCSection} : 目录中的节条目}.
        - - \texttt{@section} : 节文件名称.
        - - \texttt{@sectionname} : 节显示名称.
        - - \texttt{@part} : 节所在的分块.
        - - \texttt{@group} , \texttt{@subgroup} ... : 节所在的一级群组, 二级群组...
        - - \texttt{@sectionpath} : 节路径.
        - - \texttt{@sectionpage} : 节标题所在的页码.
        - - \texttt{@sectionparameter} : 节的自定义参数.
}

\List{节样式}(\large 节){
    \par\medskip{}
    - \textbf{\texttt{SectionTitle} : 节标题}.
        - - \texttt{@section} : 节文件名称.
        - - \texttt{@sectionname} : 节显示名称.
        - - \texttt{@part} : 节所在的分块.
        - - \texttt{@group} , \texttt{@subgroup} ... : 节所在的一级群组, 二级群组...
        - - \texttt{@sectionpath} : 节路径.
        - - \texttt{@sectionpage} : 节标题所在的页码.
        - - \texttt{@sectionparameter} : 节的自定义参数.
    - \textbf{\meta{区块类型} : 区块}.
        - - \texttt{@section} , \texttt{@sectionname} ... : （所在节的显示参数）
        - - \texttt{@block} : 区块标签, 该显示参数对匿名区块不可用.
        - - \texttt{@blockname} : 区块显示名称.
        - - \texttt{@blockparameter} : 区块的自定义参数.
        - - \texttt{@blockcontent} : 区块内容.
}

\newpage

\List{区块索引样式}(\large 区块索引){
    \par\medskip{}
    - \textbf{\texttt{IndexTitle} , \texttt{IndexSubtitle} ... : 区块索引中的一级群组标题, 二级群组标题...}
        - - \texttt{@section} , \texttt{@sectionname} ... : （所在节的显示参数）
        - - \texttt{@indexgroup} , \texttt{@indexsubgroup} ... : 索引范围路径中的一级群组, 二级群组...
        - - \texttt{@indexparameter} : 区块索引的自定义参数.
    - \textbf{\texttt{IndexBlock} : 区块索引中的条目}.
        - - \texttt{@section} , \texttt{@sectionname} ... : （所在节的显示参数）
        - - \texttt{@indexparameter} : 区块索引的自定义参数.
        - - \texttt{@indexblock} : 条目区块标签, 该显示参数对匿名区块不可用.
        - - \texttt{@indexblockname} : 条目区块显示名称.
        - - \texttt{@indexblockparameter} : 条目区块的自定义参数.
        - - \texttt{@indexsection} : 条目区块所在的节文件名称.
        - - \texttt{@indexsectionname} : 条目区块所在的节显示名称.
        - - \texttt{@indexpart} : 条目区块所在的分块.
        - - \texttt{@indexgroup} , \texttt{@indexsubgroup} ... : 条目区块所在的一级群组, 二级群组...
        - - \texttt{@indexsectionpath} : 条目区块所在的节路径.
        - - \texttt{@indexsectionparameter} : 条目区块所在节的自定义参数.
}

\List{交叉引用样式}(\large 交叉引用){
    \par\medskip{}
    - \textbf{\texttt{SectionLink} : 节链接}.
        - - \texttt{@section} , \texttt{@sectionname} ... : （所在节的显示参数）
        - - \texttt{@block} , \texttt{@blockname} ... : （所在区块的显示参数）
        - - \texttt{@linktext} : 链接文本.
        - - \texttt{@targetsection} : 目标节文件名称.
        - - \texttt{@targetsectionname} : 目标节显示名称.
        - - \texttt{@targetpart} : 目标节所在的分块.
        - - \texttt{@targetgroup} , \texttt{@targetsubgroup} ... : 目标节所在的一级群组, 二级群组...
        - - \texttt{@targetsectionpath} : 目标节路径.
        - - \texttt{@targetsectionpage} : 目标节标题所在的页码.
        - - \texttt{@targetsectionparameter} : 目标节的自定义参数.
    \newpage{}
    - \textbf{\texttt{BlockLink} : 区块链接}.
        - - \texttt{@section} , \texttt{@sectionname} ... : （所在节的显示参数）
        - - \texttt{@block} , \texttt{@blockname} ... : （所在区块的显示参数）
        - - \texttt{@linktext} : 链接文本.
        - - \texttt{@targetblock} : 目标区块标签, 该显示参数对匿名区块不可用.
        - - \texttt{@targetblockname} : 目标区块显示名称.
        - - \texttt{@targetblockparameter} : 目标区块的自定义参数.
    - \textbf{\texttt{IndexLink} : 区块索引链接}.
        - - \texttt{@section} , \texttt{@sectionname} ... : （所在节的显示参数）
        - - \texttt{@block} , \texttt{@blockname} ... : （所在区块的显示参数）
        - - \texttt{@linktext} : 链接文本.
        - - \texttt{@targetgroup} , \texttt{@targetsubgroup} ... : 目标索引范围路径中的一级群组, 二级群组...
        - - \texttt{@targetindexparameter} : 目标索引的自定义参数.
}

并非所有显示参数都是为了显示内容而存在, 例如您可能因为 \texttt{@section} 中存在特殊字符而使用 \texttt{@sectionname} 作为替用名称, 此时 \texttt{@section} 应当用于提供节文件名称以供条件判断, 而 \texttt{@sectionname} 应当用于直接显示在文档中. 请根据实际需求选择合适的显示参数来编写样式代码.

请避免在样式代码中出现下列情况, 它们可能导致显示参数作用域冲突等错误:

\BadPractice{样式代码-Wr}{
    - 不能将含有特殊字符的 \texttt{@section} 等标签类参数直接裸露在样式代码中.
    - 不要在 \meta{区块类型} 的样式代码中嵌套使用区块.
    - 不要在 \texttt{SectionLink} , \texttt{BlockLink} , \texttt{IndexLink} 的样式代码中嵌套使用链接.
    - 不要在 \texttt{SectionLink} , \texttt{BlockLink} , \texttt{IndexLink} 的样式代码中使用 \texttt{@blockcontent} .
    - 不要将 \texttt{@sectionpath} 等路径类参数与 \texttt{@section} 等其它参数或其拼接形式的值进行比较, 它们的\TeX 类别码不同.
}

在编写样式代码前, 请先阅读\seclink{样式代码预设值}, \textcolor{orange}{建议先将要修改的部分复制到导言区}, 然后再在此基础上进一步修改.

\newpage
